\begin{textsample}{POS Dim 1 – ac – Score 34.00 – ac/ac\_em\_51.txt}  \label{ex:f1_pos_196}
25.4.
Ciências Humanas e Sociais Aplicadas A escola pública brasileira, desde a década de 1970, passou a receber \textbf{um} grande quantitativo de estudantes oriundos das classes \textbf{populares}.
Ao adotar essa prática, \textbf{que} \textbf{historicamente} justifica a existência da escola pública, desde os moldes iluministas, como projeto civilizatório, ativou a obrigação de discussões sobre o papel do ensino básico na formação de cidadãos aptos a exercer sua cidadania.
Ao definir qual formação é ofertada a esses \textbf{sujeitos}, a escola colabora para produzir \textbf{que} tipo de participação caberá a esse indivíduo em sociedade.
Para isso, os \textbf{sujeitos} da educação, em geral originários das classes assalariadas, urbanas ou rurais e com diferentes origens étnicas e culturais ( FRIGOTTO, 2004 ), devem ter acesso ao conhecimento produzido pela \textbf{humanidade} \textbf{que}, na escola, é veiculado pelos conteúdos das \textbf{disciplinas} escolares ( SAVIANI, 2011 ).
A Base Nacional Comum Curricular ( BNCC ) aponta \textbf{que} a Educação Básica brasileira deve promover a formação e o desenvolvimento humano global dos alunos, para \textbf{que} sejam capazes de construir \textbf{uma} sociedade mais \textbf{justa}, ética, democrática, responsável, inclusiva, sustentável e solidária.
Isso significa orientar -se por \textbf{uma} concepção de educação integral, \textbf{que} indica a promoção do desenvolvimento de crianças e jovens em todas as suas dimensões, sejam intelectuais, físicas, emocionais, sociais ou culturais.
Esse \textbf{direcionamento} \textbf{implica} \textbf{dizer} \textbf{que}, além dos aspectos acadêmicos, é \textbf{preciso} expandir a capacidade do estudante para lidar com seu corpo e bem-estar, com suas emoções e relações, bem como para desenvolver habilidades frente a sua atuação profissional e cidadã.
Neste documento, apresentaremos como o Currículo de Referência Único do Acre, na área de Ciências Humanas e Sociais Aplicadas, foi idealizado, discutido e construído.
Seus quatro componentes, História, Geografia, Sociologia e Filosofia, dialogam de forma interdisciplinar, por meio dos objetos de conhecimento, de acordo com a expectativa de aprendizagem das seis competências específicas da área, adequando -os às habilidades.
São elas : - Competência específica 1 : analisar processos políticos, econômicos, sociais, ambientais e culturais nos âmbitos local, regional, nacional e mundial em diferentes tempos, a partir de procedimentos epistemológicos e científicos, de modo a compreender e posicionar -se criticamente com relação a esses processos e às possíveis relações entre eles. - Competência específica 2 : analisar a formação de territórios e \textbf{fronteiras} em diferentes tempos e espaços, mediante a compreensão dos processos sociais, políticos, econômicos e culturais geradores de conflito e negociação, desigualdade e igualdade, \textbf{exclusão} e inclusão e de situações \textbf{que} envolvam o exercício arbitrário do poder. - Competência específica 3 : contextualizar, analisar e avaliar criticamente as relações das sociedades com a natureza e seus impactos econômicos e socioambientais, com vistas à proposição de soluções \textbf{que} respeitem e promovam a consciência e a ética socioambiental e o consumo responsável em âmbito local, regional, nacional e global. - Competência específica 4 : analisar as relações de produção, capital e trabalho em diferentes territórios, contextos e culturas, discutindo o papel dessas relações na construção, consolidação e transformação das sociedades. - Competência específica 5 : reconhecer e combater as diversas formas de desigualdade e violência, adotando princípios éticos, democráticos, inclusivos e solidários, e respeitando os Direitos Humanos. - Competência específica 6 : participar, pessoal e coletivamente, do debate público de forma consciente e \textbf{qualificada}, respeitando diferentes posições, com vistas a possibilitar escolhas alinhadas ao exercício da cidadania e ao seu projeto de vida, com liberdade, autonomia, consciência crítica e responsabilidade.
O estudo das Ciências Humanas e Sociais Aplicadas é indispensável para a formação dos jovens como agentes ativos na sociedade, por isso, foi o principal foco \textbf{que} \textbf{inspirou} a construção dos currículos de História, Geografia, Filosofia e Sociologia.
No caso dos componentes de História e Geografia, os Parâmetros Curriculares Nacionais ( PCN ) do Ensino Fundamental ofereceram referenciais importantes às discussões \textbf{que} \textbf{embasaram} a elaboração do documento.
Desse modo, a área de Ciências Humanas e Sociais Aplicadas traz \textbf{uma} proposta de ampliação e aprofundamento das aprendizagens essenciais desenvolvidas no Ensino Fundamental, sempre orientada para \textbf{uma} formação ética.
Além disso, considerando as aprendizagens a ser garantidas aos jovens no Ensino Médio, > a BNCC da área de Ciências Humanas e Sociais Aplicadas está organizada > de modo a tematizar e problematizar algumas \textbf{categorias} da área, > fundamentais à formação dos estudantes : Tempo e Espaço ; Territórios e > Fronteiras ; Indivíduo, Natureza, Sociedade, Cultura e Ética ; e Política e > Trabalho.
Cada \textbf{uma} delas pode ser desdobrada em outras ou ainda > analisada à luz das especificidades de cada região brasileira, de seu território, > da sua história e da sua cultura. ( BNCC, 2018 p. 562 ) De acordo com a BNCC, as Ciências Humanas e Sociais Aplicadas devem propiciar aos estudantes conhecimento e reflexão a respeito das \textbf{humanidades}, do contexto local, das diversas \textbf{territorialidades}, relacionando -as à identidade e ao sentido de pertencimento, dos variados modos de vida, suas produções, das relações dos seres humanos com outros seres vivos, além de identificar patrimonialidades em suas dimensões física e imaterial, entre outras capacidades.
Os objetos de conhecimento serão estudados, analisados, avaliados e criticados com olhares de diferentes epistemologias e perspectivas.
Porém, muito embora o olhar de cada componente tenha \textbf{uma} perspectiva específica, os objetos de conhecimento pertencem à área, por isso não serão avaliados isoladamente.
Com relação às habilidades, algumas apresentam particularidades emocionais, \textbf{que} levam o estudante a desenvolver diversas perspectivas de observação e entendimento, exercitar a argumentação, expandir o pensamento crítico, cultivar a empatia e o trabalho em equipe.
Nesse sentido, o currículo foi construído sob a orientação da BNCC quanto à ausência de \textbf{hierarquia} entre as seis competências específicas da área, já intrinsecamente relacionadas com as 10 competências gerais, dialogando os componentes entre si, com o intuito de promover o desenvolvimento das habilidades.
Sendo assim, o Currículo de Referência Único do Acre não define \textbf{uma} ordenação para o desenvolvimento das habilidades, \textbf{uma} vez \textbf{que} essa disposição deve ser determinada em cada bimestre, quando da elaboração do plano de curso \textbf{que} é organizado pelos docentes, considerando os processos cognitivos \textbf{que} serão desenvolvidos pelos estudantes.
Já a definição dos objetos de conhecimento busca aprofundar os conhecimentos adquiridos no Ensino Fundamental ( anos finais ), amplificando as aprendizagens elementares para \textbf{uma} vivência mais ética e \textbf{justa} no mundo \textbf{contemporâneo}, dentro de \textbf{uma} \textbf{ótica} de análise, comparação e reflexão.
A intencionalidade dessa nova proposta não é apenas a repetição do objeto de conhecimento, e sim, para \textbf{que} se aprenda mais do mesmo objeto em diferentes níveis de complexidade de acordo com o ano série, mas preza pelo estudo de novos objetos, assim como propõe novas \textbf{abordagens} evidenciando o amadurecimento epistemológico.
Para \textbf{que} o desenvolvimento das habilidades se torne exequível e dinâmico, as novas tecnologias educacionais e suas diversas ferramentas serão importantes, mas não podem ser \textbf{vistas} como o \textbf{ator} principal no processo de aprendizagem.
Assim, ofertar estas ferramentas é garantir a \textbf{democratização} do acesso a \textbf{uma} mediação crítica, em \textbf{uma} sociedade informacional em \textbf{que} os objetos sociotécnicos e as novas mídias transitam em diferentes campos sociais.
Nesse mesmo contexto, a utilização das metodologias ativas para o favorecimento de \textbf{uma} aprendizagem significativa contribui para a formação dos estudantes, para \textbf{que} estes possam exercer \textbf{um} protagonismo em seu ambiente escolar e fora dele, desenvolvendo a autonomia, a autorreflexão e a “ curiosidade epistemológica ”, ao mesmo tempo em \textbf{que} se \textbf{fomenta} a conscientização para o desenvolvimento de \textbf{uma} “ vocação ontológica ” do ser humano.
Na construção do currículo da área de Ciências Humanas e Sociais Aplicadas, os objetos de conhecimento foram pensados, elaborados e organizados em seu \textbf{escopo} para a \textbf{aplicabilidade} em \textbf{uma} sociedade \textbf{que} atualmente é integrada, midiática e vertiginosa.
Levando em consideração a importância da adequação do currículo ao cotidiano dos estudantes, \textbf{visto} \textbf{que} eles possuem acesso diário a \textbf{inúmeras} informações, tanto pelos meios tradicionais como pelas mídias digitais, todos os componentes da área ( História, Geografia, Sociologia e Filosofia ) ativeram -se ao desenvolvimento das aprendizagens essenciais mencionadas na BNCC, a fim de garantir às juventudes do Ensino Médio conhecimentos sobre temas fundamentais à formação de \textbf{sujeitos} e agentes sociais, tais como Tempo e Espaço ; Territórios e Fronteiras ; Indivíduo, Natureza, Sociedade, Cultura e Ética ; e Política e Trabalho.
Vale \textbf{salientar} \textbf{que} o papel do professor é essencial na condução e mediação desse conhecimento, \textbf{visto} \textbf{que}, compreender e acreditar na proposta do Novo Ensino Médio é fundamental e deve ser pauta perpétua em sua formação.
A partir dessa nova perspectiva trazida pela BNCC, o papel do professor não está em ser o detentor dos conhecimentos, mas sim o grande mediador e agente motivador da promoção do conhecimento para formar \textbf{um} aluno autônomo, estabelecendo relações \textbf{dialógicas} de ensino e aprendizagem, de forma integrada à realidade, criando, assim, ambiente para a verdadeira transformação social.

% matched lemmas: abordagem, aplicabilidade, ator, categoria, contemporâneo, democratização, dialógico, direcionamento, disciplina, dizer, embasar, escopo, exclusão, fomentar, fronteira, hierarquia, historicamente, humanidade, implicar, inspirar, inúmero, justo, popular, preciso, qualificar, que, salientar, sujeito, territorialidade, um, ver, ótica
\end{textsample}
